\chapter{Portfolio shapes that are reported to work}
\label{ch:portfolio}

\section*{What this chapter is for}

Four shapes, what each one proves, what each costs, and which to build first.
The ordering is not arbitrary: one of the four is repeatedly singled out in
reports of what separates candidates, and it is not the one most people build.

\section{The eval harness}

\textbf{What it is.} A small system that decides, mechanically, whether a
model-driven behaviour is still correct, and that can say what it would take to
make it fail.

\textbf{Why it is first.} It is the differentiating artefact. \reported{Some
companies ask for portfolio evidence of a retrieval system \emph{with} an
evaluation harness specifically to separate people who have built something
from people who have followed a
tutorial.}{field-guide-get-hired} In take-home
assessments, \reported{not starting with evaluation is reported as a red
flag.}{field-guide-take-home}

\textbf{What it proves.} That you have been hurt by a non-deterministic system
and did something structural about it. That is the domain doubt and the
judgement doubt closed at once.

\textbf{What it costs.} A week, if the thing being evaluated already exists.
Much less if you evaluate something small.

\begin{kitnote}
It does not have to be large. A harness with a dozen scenarios, a fixed set of
assertions, a recorded baseline, and two deliberately broken variants that the
suite catches is enough to have the conversation. Chapter~\ref{ch:evals} is the
content; this is about the fact that it exists and is yours.
\end{kitnote}

\section{Cost arithmetic, done out loud}

\textbf{What it is.} A written analysis of what a system costs per request and
in aggregate, with the arithmetic on the page, and \dash{} if you have it
\dash{} a before-and-after from a change you made.

\textbf{Why it works.} Cost and latency are \reported{among the most frequently
asked topics}{field-guide}, and almost nobody arrives able to do the
multiplication in front of an interviewer. A written version means you have
done it once slowly.

\textbf{What it proves.} That you think about the thing that decides whether a
feature ships. \judgement{It is also the single clearest way to sound like
somebody who has operated a system rather than built one, which is most of what
Staff assessment is looking for.}

\textbf{What it costs.} An afternoon. It is the cheapest item in this chapter by
a wide margin and the most under-used.

\section{A written opinion}

\textbf{What it is.} One or two pages taking a position on something contested
in this field, with the reasoning and the conditions under which you would
change your mind.

\textbf{Why it works.} \reported{At least one employer asks candidates directly
for their view on where organisations go wrong with this technology, in the form
of a short essay.}{field-guide-get-hired} More generally, judgement is what is
being hired at this level and an essay is the only artefact that demonstrates it
directly rather than as a by-product.

\textbf{What it proves.} That you have a position, that you can hold it under
scrutiny, and \dash{} if you write the last paragraph properly \dash{} that you
know what would falsify it.

\textbf{What it costs.} A day. \judgement{It is also the artefact most likely to
be read, because it is the only one that can be read on a phone.}

\section{A deployed integration}

\textbf{What it is.} A working system, at a URL, that does one real thing with
a model and a tool.

\textbf{Why it works.} \reported{Deployed and clickable beats a
notebook}{field-guide-get-hired}, and a URL answers the question of whether it
runs anywhere other than your machine without anybody having to ask it.

\textbf{What it proves.} Currency and the ability to finish. Both matter for a
candidate whose job history is in another domain.

\textbf{What it costs.} Two or three days, most of it deployment rather than
the interesting part. Chapter~\ref{ch:surface} is about choosing \emph{what} to
integrate with, which is the decision that makes this shape worth more than the
effort it takes.

\section{Choosing}

\begin{kittable}{@{}lllX@{}}
\textbf{Shape} & \textbf{Cost} & \textbf{Closes} & \textbf{Build it when} \\
\midrule
Eval harness & a week & domain, judgement & always; build it first \\
Cost analysis & an afternoon & operating instinct & always; it is nearly free \\
Written opinion & a day & judgement, communication & you have a position worth defending \\
Deployed integration & 2--3 days & currency, finishing & you have nothing running anywhere \\
\end{kittable}

\judgement{If you have one week in total, build a small deployed thing and put
an evaluation harness on it, then spend the last afternoon on the cost
arithmetic for what you built. That combination closes every doubt an artefact
can close, and it is one artefact rather than three.}

\begin{kittrap}
The shape most people build first is a retrieval demonstration with a chat
interface. It is the most tutorialised thing in the field, which means a reader
has seen thirty of them and cannot distinguish yours from a weekend's work
following a guide. The demonstration is not the problem; the absence of
evaluation, cost analysis or a stated position is. The same repository with a
harness and a register attached is a completely different document.
\end{kittrap}

\begin{drill}
Write the one-sentence thesis for each of the four shapes as you would build
them, before building anything. If a thesis is ``it demonstrates that I can use
a vector database'', that shape is not worth your week. Keep going until at
least two of the four produce a sentence somebody could disagree with.
\end{drill}

\section*{What this chapter does not claim}

The reports behind the first and fourth shapes describe what some employers ask
for and what some candidates found worked. They do not establish how common
either is, and single striking outcomes \dash{} offers arriving within days of
publishing something \dash{} are individual anecdotes and are refused as
evidence in Appendix~\ref{app:refused}.

The ordering of the four, and the recommendation about what to do with one
week, are the author's judgement rather than a measured result.
